\documentclass[11pt,a4paper]{beamer}
\usetheme{Madrid}
\usepackage[utf8]{inputenc}
\usepackage{amsmath}
\usepackage{amsfonts}
\usepackage{amssymb}
\usepackage{graphicx}
\usepackage{xcolor}
\usepackage{tikz}
\usepackage{booktabs}
\usepackage{array}
\usepackage{colortbl}
\usepackage{multicol}

% ============================================================
% TINY FONT FOR MAXIMUM CONTENT
% ============================================================
\usefonttheme{professionalfonts}

\setbeamerfont{normal text}{size=\scriptsize}
\setbeamerfont{itemize/enumerate body}{size=\scriptsize}
\setbeamerfont{itemize/enumerate subbody}{size=\tiny}
\setbeamerfont{block body}{size=\scriptsize}
\setbeamerfont{block title}{size=\scriptsize}
\setbeamerfont{frametitle}{size=\small}
\setbeamerfont{title}{size=\normalsize}
\setbeamerfont{subtitle}{size=\small}
\setbeamerfont{author}{size=\tiny}
\setbeamerfont{date}{size=\tiny}

\renewcommand{\scriptsize}{\fontsize{6pt}{7pt}\selectfont}
\renewcommand{\footnotesize}{\fontsize{6.5pt}{8pt}\selectfont}
\renewcommand{\small}{\fontsize{7.5pt}{9pt}\selectfont}
\renewcommand{\normalsize}{\fontsize{8.5pt}{10.5pt}\selectfont}

\newcommand{\redflag}[1]{\textcolor{red}{\textbf{[!] #1}}}
\newcommand{\bluenote}[1]{\textcolor{blue}{\textbf{[i] #1}}}
\newcommand{\greennote}[1]{\textcolor{green!50!black}{\textbf{[CORRECT] #1}}}
\newcommand{\examfocus}[1]{\textcolor{orange!80!black}{\textbf{[FOCUS] #1}}}
\newcommand{\trap}[1]{\textcolor{purple}{\textbf{[TRAP] #1}}}
\newcommand{\correct}[1]{\textcolor{green!50!black}{\textbf{[right] #1}}}
\newcommand{\scenario}[1]{\textcolor{blue!70!black}{\textbf{[SCENARIO] #1}}}
\newcommand{\keyterm}[1]{\textcolor{red!80!black}{\textbf{[KEY] #1}}}

\title[CAMS Domain C - Missing Hard Questions]{CAMS Exam: Domain C Missing Content\\25 Hard Questions Missing from Standard Materials}
\author{Based on 2024-2026 Candidate Feedback}
\date{}

\setbeamercovered{transparent}
\setbeamertemplate{navigation symbols}{}
\setbeamertemplate{frametitle continuation}{}
\setlength{\parskip}{2pt}
\setlength{\itemsep}{1pt}

\begin{document}

% ============================================================
% TITLE SLIDE
% ============================================================
\begin{frame}
\titlepage
\centering
\tiny Domain C: 30\% of Exam | Risk Assessment Methodology, Board Minutes, CCO Independence, SAR Timelines, 314(a)/314(b), Recordkeeping, and More
\end{frame}

% ============================================================
% SECTION 1: RISK ASSESSMENT METHODOLOGY (MISSING DETAIL)
% ============================================================

\begin{frame}{Risk Assessment - Inherent vs. Residual Risk Calculation}
\begin{block}{Core Concept - Heavily Tested}
Inherent risk = risk before controls. Residual risk = risk after controls.
\end{block}
\textbf{Calculation Method:}
\begin{itemize}
\item Inherent risk scored 0-100 (based on products, customers, geography, channels)
\item Control effectiveness scored 0-3 (0=no control, 1=weak, 2=moderate, 3=strong)
\item Residual risk = Inherent risk ÷ Control effectiveness factor
\item Or: Residual risk = Inherent risk - (Inherent risk multiply Control effectiveness percentage)
\end{itemize}
\scenario{Bank calculates inherent risk for MSB customer = 80/100. Control effectiveness = 2.5/3. What is residual risk?}
\trap{Subtract control effectiveness from inherent risk: 80 - 2.5 = 77.5}
\correct{Residual risk = Inherent risk ÷ (Control effectiveness as factor). 80 dividde 2.5 = 32. Residual risk is significantly reduced by strong controls.}
\redflag{High inherent risk is acceptable if residual risk is within appetite.}
\end{frame}

\begin{frame}{Risk Appetite Statement - What It Must Include}
\begin{block}{Board-Approved Risk Appetite}
Risk appetite defines acceptable residual risk levels.
\end{block}
\textbf{Mandatory Components:}
\begin{itemize}
\item Quantitative thresholds (e.g., maximum residual risk score of 40 for any relationship)
\item Qualitative boundaries (e.g., no correspondent relationships with shell banks)
\item Escalation triggers (e.g., any relationship exceeding appetite requires Board approval)
\item Review frequency (at least annually)
\end{itemize}
\scenario{Bank has no written risk appetite statement. CCO argues risk assessment alone is sufficient.}
\trap{Risk assessment replaces need for separate risk appetite statement.}
\correct{Risk appetite statement is required. Board must define acceptable risk levels.}
\end{frame}

\begin{frame}{Dynamic Risk Reassessment - Trigger Events}
\begin{block}{Risk is NOT Static}
Customer risk ratings must be updated when trigger events occur.
\end{block}
\textbf{Mandatory Reassessment Triggers:}
\begin{itemize}
\item SAR filing on the customer
\item Negative media or law enforcement inquiry
\item Significant change in transaction volume or pattern
\item Customer begins using new high-risk products
\item Change in beneficial ownership
\item Customer moves to/from high-risk jurisdiction
\item Regulatory finding related to customer type
\item Merger or acquisition involving customer
\end{itemize}
\scenario{Bank reviews customer risk ratings annually only. Customer files SAR but rating remains unchanged for 11 months until annual review.}
\trap{Annual review is sufficient for all customers regardless of events.}
\correct{SAR filing is a trigger requiring IMMEDIATE reassessment. Cannot wait for annual review.}
\end{frame}

% ============================================================
% SECTION 2: BOARD MEETING MINUTES (MISSING FORMAT)
% ============================================================

\begin{frame}{Board Minutes - Required Content for AML Approval}
\begin{block}{What Regulators Expect in Minutes}
Minutes must demonstrate meaningful oversight, not rubber-stamping.
\end{block}
\textbf{Required Elements:}
\begin{itemize}
\item Who attended (names and titles)
\item Discussion summary (risks, metrics, challenges)
\item Questions asked by Board members
\item Management responses
\item Any dissents or concerns
\item Vote results (unanimous or otherwise)
\item Specific approval of AML policy (not just general)
\item Date of next AML review
\end{itemize}
\scenario{Board minutes state: AML policy approved unanimously. No other text. CCO argues approval is valid.}
\trap{Verbal approval with minutes entry is sufficient regardless of detail.}
\correct{Minutes must show evidence of meaningful discussion. One-line entry suggests rubber-stamping and is a regulatory deficiency.}
\end{frame}

\begin{frame}{Board Minutes - Missing Discussion Red Flag}
\scenario{Regulator reviews Board minutes for past 3 years. Each quarter: AML policy reviewed and approved. No mention of SAR volumes, backlog, audit findings, or regulatory changes. No questions from Board recorded.}
\vspace{0.2cm}
\textbf{What the Exam Tests:} Insufficient Board oversight.
\vspace{0.2cm}
\textbf{Analysis:}
\begin{enumerate}
\item Repetitive one-line approval suggests lack of engagement
\item No discussion of metrics = no oversight
\item Board cannot approve policy without understanding risks
\item Regulator will find governance deficiency
\end{enumerate}
\trap{Minutes show approval occurred, which satisfies regulatory requirements.}
\correct{Minutes must document substantive discussion. Lack of detail is evidence of Board non-engagement.}
\end{frame}

% ============================================================
% SECTION 3: CCO INDEPENDENCE (EXPANDED)
% ============================================================

\begin{frame}{CCO Independence - Reporting Line Requirements}
\begin{block}{CCO Cannot Report to Revenue Functions}
The BSA Officer must have direct Board access without filtering.
\end{block}
\textbf{Acceptable Reporting Lines:}
\begin{itemize}
\item Direct to Board of Directors
\item Direct to Audit Committee of the Board
\item Administrative reporting to CEO (but with direct Board access)
\end{itemize}
\textbf{UNACCEPTABLE Reporting Lines:}
\begin{itemize}
\item Reporting to Head of Retail Banking (revenue)
\item Reporting to Head of International Banking (revenue)
\item Reporting to Chief Financial Officer (may have revenue conflicts)
\item Any line where CCO's budget is controlled by business lines
\end{itemize}
\redflag{If CCO can be overridden by business lines for commercial reasons, CCO lacks required independence.}
\end{frame}

\begin{frame}{CCO Independence Violation - Revenue Pressure}
\scenario{CCO reports to Head of Corporate Banking. Head instructs CCO to delay SAR filing on \$100M revenue client. CCO objects but Head says: I'm your boss, do it. CCO delays filing.}
\vspace{0.2cm}
\textbf{What the Exam Tests:} CCO reporting line independence.
\vspace{0.2cm}
\textbf{Analysis:}
\begin{enumerate}
\item CCO should report to Board or Audit Committee, not revenue function
\item Head of Corporate Banking has inherent revenue conflict
\item Delaying SAR for commercial reasons is violation
\item CCO should escalate to Board/regulator despite reporting line
\end{enumerate}
\trap{CCO followed management direction; no personal violation.}
\correct{CCO has independent regulatory obligation. Reporting line does not excuse violation. CCO must file SAR regardless of management instruction.}
\end{frame}

\begin{frame}{CCO Resources - Adequate Staffing Requirement}
\begin{block}{Under-Resourcing is NOT a Defense}
Regulators will cite inadequate resources as a violation.
\end{block}
\textbf{Evidence of Resource Inadequacy:}
\begin{itemize}
\item Persistent backlog of uninvestigated alerts (>30 days aging)
\item High analyst turnover (burnout, underpaid)
\item Manual processes for high-volume tasks
\item Delayed SAR filings due to insufficient reviewers
\item Board rejects reasonable budget requests from CCO
\item Staff lack required qualifications
\end{itemize}
\scenario{Bank has 40,000 alert backlog. CCO requested 10 additional analysts. Board approved 2. CCO argues insufficient resources caused backlog.}
\trap{Board controls budget; CCO cannot be held responsible.}
\correct{CCO must document resource requests and escalate. Bank remains responsible for backlog. Regulator will cite inadequate resources.}
\end{frame}

% ============================================================
% SECTION 4: REGULATORY REPORTING DEADLINES (MISSING)
% ============================================================

\begin{frame}{SAR Filing Deadlines - 30/60 Day Rule}
\begin{block}{SAR Deadline Depends on Suspect Identification}
Different deadlines apply based on whether suspect is identified.
\end{block}
\begin{tabular}{p{0.45\textwidth}|p{0.45\textwidth}}
\hline
\rowcolor{lightgray}
\textbf{Situation} & \textbf{Deadline} \\
\hline
Suspect identified & 30 calendar days from suspicion \\
Suspect NOT identified & 60 calendar days from suspicion \\
\hline
\end{tabular}
\vspace{0.2cm}
\textbf{Key Points:}
\begin{itemize}
\item Clock starts when suspicion is reached, not when investigation is complete
\item Do NOT delay filing to identify suspect
\item File by day 30 if suspect known, day 60 if unknown
\item If suspect identified after filing, no need to amend
\end{itemize}
\scenario{Analyst reaches suspicion on Day 1 but cannot identify suspect. Suspect identified on Day 45. When must SAR be filed?}
\trap{SAR must be filed within 30 days of suspect identification (Day 75).}
\correct{SAR must be filed within 60 days of suspicion (Day 60). Suspect identification does NOT restart clock.}
\end{frame}

\begin{frame}{CTR Filing Deadline - 15 Day Rule}
\begin{block}{Currency Transaction Report Deadline}
CTR must be filed within 15 days of the reportable transaction.
\end{block}
\textbf{Key Rules:}
\begin{itemize}
\item Single transaction exceeding \$10,000 in cash
\item Multiple related transactions exceeding \$10,000 on same business day
\item Aggregation required for same customer on same day
\item Filed with FinCEN via BSA E-Filing system
\end{itemize}
\scenario{Customer makes \$9,500 cash deposit at Branch A and \$9,500 cash deposit at Branch B on same day. Each branch files CTR?}
\trap{Each branch files separate CTR for \$9,500 (below threshold).}
\correct{Bank must aggregate related transactions. Combined \$19,000 exceeds threshold. Single CTR required.}
\end{frame}

\begin{frame}{314(a) Response Deadline - 14 Business Days}
\begin{block}{FinCEN 314(a) - Mandatory Law Enforcement Requests}
Financial institutions MUST respond within 14 business days.
\end{block}
\textbf{314(a) Process:}
\begin{enumerate}
\item FinCEN issues request with subject names
\item Institution searches accounts and transactions
\item Positive matches reported within 14 business days
\item Negative responses NOT required (but best practice)
\item Maintain confidentiality - cannot notify customer
\end{enumerate}
\scenario{Bank receives 314(a) request on Friday. Compliance manager says: We have 14 calendar days, so due in 2 weeks.}
\trap{14 calendar days from receipt is correct.}
\correct{314(a) deadline is 14 BUSINESS days, not calendar days. Weekends and holidays excluded.}
\end{frame}

\begin{frame}{FBAR Filing Deadline - April 15}
\begin{block}{Foreign Bank Account Report (FBAR)}
US persons with foreign accounts exceeding \$10,000 aggregate must file FBAR.
\end{block}
\textbf{Deadlines:}
\begin{itemize}
\item Annual filing due April 15
\item Automatic extension to October 15 (no form required)
\item No further extensions available
\item Electronic filing only (FinCEN Form 114)
\end{itemize}
\scenario{US citizen has \$15,000 in foreign account. Misses April 15 deadline. Files on May 1.}
\trap{Late filing is acceptable as long as filed within 6 months.}
\correct{FBAR has automatic extension to October 15. May 1 filing is acceptable but late. No penalty for reasonable cause.}
\end{frame}

% ============================================================
% SECTION 5: USA PATRIOT ACT DETAILS (MISSING)
% ============================================================

\begin{frame}{Section 311 - Primary Money Laundering Concern}
\begin{block}{FinCEN Can Designate Foreign Institutions}
Three escalating levels of special measures.
\end{block}
\textbf{Level 1 - Special Due Diligence:}
\begin{itemize}
\item Enhanced recordkeeping
\item Additional reporting requirements
\item Information collection
\end{itemize}
\textbf{Level 2 - Enhanced Due Diligence:}
\begin{itemize}
\item Identify beneficial owners
\item Monitor transactions
\item Restrict account types
\end{itemize}
\textbf{Level 3 - Prohibition:}
\begin{itemize}
\item U.S. banks cannot open or maintain correspondent accounts
\item Existing accounts must be closed
\end{itemize}
\redflag{Section 311 applies to FOREIGN institutions, not U.S. banks directly.}
\end{frame}

\begin{frame}{Section 313 - Shell Bank Prohibition}
\begin{block}{Definition of Shell Bank}
A bank with NO physical presence in any country and NOT affiliated with a regulated financial institution.
\end{block}
\textbf{Absolute Prohibition:}
\begin{itemize}
\item U.S. banks CANNOT maintain correspondent accounts for foreign shell banks
\item NO exceptions - not for any reason
\item Violation results in penalties and potential termination orders
\end{itemize}
\textbf{NOT a Shell Bank:}
\begin{itemize}
\item Bank with physical office and staff in home country
\item Bank affiliated with regulated depository institution
\item Bank with license and physical presence
\end{itemize}
\scenario{U.S. bank discovers respondent bank has no physical office, operates from shared workspace, no employees.}
\trap{Bank has license, so not a shell bank.}
\correct{NO physical presence + NOT affiliated = shell bank. Account must be terminated immediately.}
\end{frame}

\begin{frame}{Section 319(b) - Foreign Bank Record Production}
\begin{block}{Records Must Be Available in US}
Foreign banks with correspondent accounts must maintain records in US or designate agent.
\end{block}
\textbf{Key Requirements:}
\begin{itemize}
\item Foreign bank must designate agent for service of process in US
\item Must maintain records related to correspondent account
\item Records must be produced within specified timeframe
\item Failure to produce = grounds for account termination
\item Funds may be subject to forfeiture
\end{itemize}
\scenario{US subpoena requests records from foreign respondent bank. Foreign bank refuses.}
\trap{US bank has no obligation since foreign bank is separate entity.}
\correct{US bank may be required to terminate correspondent relationship. Funds in account may be seized under 319(a).}
\end{frame}

% ============================================================
% SECTION 6: 314(a) AND 314(b) DETAILS (MISSING)
% ============================================================

\begin{frame}{314(a) vs. 314(b) - Key Differences}
\begin{block}{Two Different Information Sharing Mechanisms}
Both serve different purposes with different requirements.
\end{block}
\begin{tabular}{p{0.35\textwidth}|p{0.30\textwidth}|p{0.30\textwidth}}
\hline
\rowcolor{lightgray}
\textbf{Feature} & \textbf{314(a)} & \textbf{314(b)} \\
\hline
Direction & LE to FI & FI to FI \\
Voluntary? & MANDATORY & Voluntary (opt-in) \\
Trigger & FinCEN request & Institution initiative \\
Response deadline & 14 business days & N/A \\
Safe harbor & Yes & Yes \\
Tipping-off? & Cannot notify & Not tipping-off \\
\hline
\end{tabular}
\vspace{0.2cm}
\redflag{314(a) is MANDATORY. 314(b) is VOLUNTARY but requires opt-in.}
\end{frame}

\begin{frame}{314(b) Opt-In Requirement - Missing Detail}
\begin{block}{Must File Notice Before Sharing}
Institution cannot share 314(b) information without opt-in.
\end{block}
\textbf{Opt-In Process:}
\begin{enumerate}
\item File FinCEN Form 101 (one-time)
\item Designate points of contact
\item Maintain confidentiality of shared information
\item Can share with other opted-in institutions only
\item Retroactive sharing NOT permitted
\end{enumerate}
\scenario{Bank never filed 314(b) opt-in. Shares SAR information with another bank to investigate common customer.}
\trap{Sharing SAR information is permitted under safe harbor regardless of opt-in.}
\correct{314(b) opt-in is REQUIRED before sharing. Sharing without opt-in violates tipping-off prohibition.}
\end{frame}

\begin{frame}{314(b) - What CAN Be Shared}
\begin{block}{Permissible Information Sharing}
Only information related to suspected ML/TF activity.
\end{block}
\textbf{Can Share:}
\begin{itemize}
\item Suspicious transaction patterns
\item Identities of suspected persons or entities
\item Account information (with safeguards)
\item SARs (with other opted-in institutions)
\item Typologies and red flags
\end{itemize}
\textbf{Cannot Share:}
\begin{itemize}
\item Information for non-AML purposes (marketing, competitive intelligence)
\item With institutions not opted into 314(b)
\item With public or media
\item With customers or subjects of investigation
\item Confidential customer information unrelated to suspicion
\end{itemize}
\redflag{314(b) sharing is NOT a license to share all customer information. Must be suspicion-based.}
\end{frame}

% ============================================================
% SECTION 7: THIRD-PARTY RISK (EXPANDED)
% ============================================================

\begin{frame}{Vendor Due Diligence - Pre-Engagement Requirements}
\begin{block}{Due Diligence BEFORE Contract Signing}
Cannot outsource AML functions without proper vendor vetting.
\end{block}
\textbf{Required Pre-Engagement Diligence:}
\begin{itemize}
\item Vendor background and reputation checks
\item Financial stability review
\item Vendor's own AML program assessment
\item Data security and breach history
\item Subcontractor relationships
\item References from other bank clients
\item Regulatory enforcement history
\item Business continuity and disaster recovery
\item Right-to-audit provisions
\end{itemize}
\scenario{Bank signs contract with AML monitoring vendor. Conducts due diligence 6 months AFTER implementation.}
\trap{Due diligence can be performed post-implementation as long as vendor is reputable.}
\correct{Due diligence MUST be performed BEFORE contract execution. Post-implementation diligence is a violation.}
\end{frame}

\begin{frame}{Vendor Oversight - Ongoing Monitoring Requirements}
\begin{block}{Due Diligence is NOT One-Time}
Ongoing monitoring required for all AML vendors.
\end{block}
\textbf{Ongoing Oversight Requirements:}
\begin{itemize}
\item Annual vendor risk reassessment
\item Service level agreement monitoring (uptime, accuracy)
\item Regular performance reviews
\item Periodic on-site audits (frequency based on risk)
\item Vendor security assessment updates
\item Subcontractor change notification and approval
\item Escalation procedures for vendor failures
\end{itemize}
\scenario{Bank conducted vendor due diligence at onboarding 5 years ago. No reviews since. Vendor performance degraded.}
\trap{Initial due diligence is sufficient; vendor responsible for performance.}
\correct{Ongoing vendor oversight is required. 5 years without review is material deficiency.}
\end{frame}

\begin{frame}{Right-to-Audit Clause - Critical Contract Provision}
\begin{block}{Bank Must Have Access to Vendor Controls}
Right-to-audit clauses enable verification of vendor compliance.
\end{block}
\textbf{What Right-to-Audit Must Include:}
\begin{itemize}
\item On-site audit access (not just remote)
\item Access to vendor's AML controls and testing results
\item Access to vendor's subcontractor information
\item Review of vendor's audit reports
\item Copy of vendor's independent testing results
\item Access to vendor's employee training records
\item Right to bring third-party auditors
\end{itemize}
\scenario{Vendor contract states: Bank may request audit results at vendor's discretion. Vendor refuses bank's audit request.}
\trap{Bank can rely on vendor's internal audit results.}
\correct{Right-to-audit must be contractual and enforceable. Vendor discretion clause is insufficient.}
\end{frame}

% ============================================================
% SECTION 8: INDEPENDENT TESTING (EXPANDED)
% ============================================================

\begin{frame}{Independent Testing - Who Can Test and Who Cannot}
\begin{block}{Independence is About the PERSON, Not the Firm}
External firm does NOT guarantee independence.
\end{block}
\begin{tabular}{p{0.45\textwidth}|p{0.45\textwidth}}
\hline
\rowcolor{lightgray}
\textbf{Acceptable Testers} & \textbf{NOT Acceptable Testers} \\
\hline
Internal Audit (if independent) & Compliance testing itself \\
External Audit Firm & Audit that designed controls \\
Third-Party Consultant & Former compliance staff (<2 years separation) \\
Independent consultant & Business line employees \\
No prior involvement & Vendors who built the system \\
New audit team members & CCO's direct reports \\
\hline
\end{tabular}
\redflag{Previous employment in compliance creates independence gap. Minimum separation period typically 1-2 years.}
\end{frame}

\begin{frame}{Testing Scope - Operational Effectiveness vs. Design}
\begin{block}{Testing Must Validate OPERATIONAL Effectiveness}
Policies on paper are insufficient; must test actual operation.
\end{block}
\textbf{Design Testing (Insufficient Alone):}
\begin{itemize}
\item Review policies and procedures
\item Confirm documentation exists
\item Verify Board approval
\end{itemize}
\textbf{Operational Testing (Required):}
\begin{itemize}
\item Test thresholds using historical SARs (back-testing)
\item Calculate false negative rate
\item Review sample of closed alerts for quality
\item Test actual response times
\item Validate data completeness (reconciliation)
\item Observe employee behavior
\end{itemize}
\scenario{Audit reviews AML policies, confirms Board approval, and checks that training records exist. No transaction testing performed.}
\trap{Policy review plus training review is sufficient for independent testing.}
\correct{Operational testing is REQUIRED. Policy-only testing is insufficient.}
\end{frame}

\begin{frame}{Audit Finding Remediation - Tracking and Validation}
\begin{block}{Findings Must Be Tracked to Closure}
Unremediated findings are ongoing deficiencies.
\end{block}
\textbf{Remediation Requirements:}
\begin{itemize}
\item Written remediation plan for each finding
\item Responsible party and due date
\item Regular status updates to management
\item Validation testing after remediation
\item Board reporting for material findings
\item Tracking system for open findings
\end{itemize}
\scenario{Audit finds 14 deficiencies. Bank remediates 12. Audit closes all findings because 12 are fixed.}
\trap{Findings are closed when remediation plan is complete.}
\correct{Each finding must be individually remediated AND validated. Cannot close findings until all are fixed.}
\end{frame}

% ============================================================
% SECTION 9: RECORDKEEPING (EXPANDED)
% ============================================================

\begin{frame}{Recordkeeping - 5-Year Retention Rule Details}
\begin{block}{Five Years From SPECIFIC Dates}
Different records have different retention start dates.
\end{block}
\begin{tabular}{p{0.40\textwidth}|p{0.55\textwidth}}
\hline
\rowcolor{lightgray}
\textbf{Record Type} & \textbf{5 Years From} \\
\hline
Account opening records & Account closing date \\
Transaction records & Transaction date \\
SAR filing & SAR filing date \\
CTR filing & CTR filing date \\
Investigation records (no SAR) & Investigation completion date \\
Training records & Training completion date \\
Audit reports & Report issuance date \\
Board minutes & Meeting date (permanent) \\
\hline
\end{tabular}
\scenario{Customer closes account. Bank deletes KYC records after 5 years from account opening.}
\trap{5 years from account opening meets requirement.}
\correct{5 years from ACCOUNT CLOSING, not opening. Deletion at 5 years from opening is premature violation.}
\end{frame}

\begin{frame}{Record Destruction - Legal Hold Suspension}
\begin{block}{Subpoena Triggers Legal Hold}
Normal record destruction must be suspended upon subpoena receipt.
\end{block}
\textbf{Legal Hold Requirements:}
\begin{itemize}
\item Identify all relevant records
\item Notify all departments (compliance, IT, records management)
\item Suspend any scheduled destruction
\item Preserve records in original format
\item Maintain chain of custody
\item Document legal hold implementation
\item Release only after authorized by legal counsel
\end{itemize}
\scenario{Bank receives subpoena for customer records. Automated system deletes records 3 days later under normal retention policy.}
\trap{Bank had no obligation to preserve since subpoena didn't specify retention.}
\correct{Legal hold is automatic upon subpoena receipt. Automated deletion after subpoena is spoliation.}
\end{frame}

% ============================================================
% SECTION 10: ESCALATION AND OFFBOARDING (EXPANDED)
% ============================================================

\begin{frame}{Escalation Matrix - Who Approves What}
\begin{block}{Different Decisions Require Different Approvals}
Escalation level depends on decision type and risk level.
\end{block}
\begin{tabular}{p{0.40\textwidth}|p{0.55\textwidth}}
\hline
\rowcolor{lightgray}
\textbf{Decision} & \textbf{Approval Required} \\
\hline
Close alert as false positive & Analyst + supervisor \\
File SAR & Analyst + supervisor \\
Maintain high-risk relationship & CCO + senior management \\
Offboard customer & Compliance committee + senior management \\
Material policy change & Board of Directors \\
Risk appetite change & Board of Directors \\
Terminate correspondent relationship & CCO + senior management \\
\hline
\end{tabular}
\redflag{CCO alone cannot approve material policy changes. Board approval required.}
\end{frame}

\begin{frame}{Offboarding - When Law Enforcement Requests Account Stay Open}
\begin{block}{Keep Open Requests Must Be Documented}
Bank should generally comply with written LE requests.
\end{block}
\textbf{Protocol for Keep Open Requests:}
\begin{enumerate}
\item Obtain request IN WRITING from law enforcement
\item Document request details (date, agency, officer)
\item Obtain senior management approval to continue
\item Maintain heightened monitoring during extension
\item Do NOT disclose request to customer
\item Document decision to keep open or close
\item Coordinate closure timing with LE
\end{enumerate}
\scenario{FBI agent orally asks bank to keep account open. Bank complies. No written documentation.}
\trap{Oral request is sufficient; documentation is administrative.}
\correct{Keep open requests must be IN WRITING. Bank should request written confirmation before agreeing.}
\end{frame}

\begin{frame}{Offboarding - Customer Funds Return Requirements}
\begin{block}{Cannot Withhold Customer Funds}
Even when offboarding for AML reasons, customer funds must be returned.
\end{block}
\textbf{Funds Return Requirements:}
\begin{itemize}
\item Return funds promptly after account closure
\item Provide check or wire to verified account
\item Cannot delay return as penalty
\item Cannot condition return on customer signing waiver
\item Cannot donate unclaimed funds to charity
\item Escheat unclaimed funds to state after required period
\end{itemize}
\scenario{Bank closes account for money laundering suspicion. Withholds \$50,000 balance pending investigation. 6 months later, still holding funds.}
\trap{Bank can hold funds indefinitely while investigating.}
\correct{Bank must return funds promptly unless LE requests freeze. Withholding without legal authority is conversion.}
\end{frame}

% ============================================================
% SECTION 11: EMPLOYEE DUE DILIGENCE (EXPANDED)
% ============================================================

\begin{frame}{Background Checks - Statutory Disqualification}
\begin{block}{Certain Convictions Prohibit Employment}
Persons convicted of money laundering or dishonesty crimes may be prohibited.
\end{block}
\textbf{Section 19 of FDI Act (US):}
\begin{itemize}
\item Persons convicted of crimes involving dishonesty, breach of trust, or money laundering
\item Cannot become institution-affiliated parties without prior written consent of FDIC
\item Applies to directors, officers, employees, agents
\item Consent required even after sentence completion
\end{itemize}
\scenario{Bank hires individual convicted of fraud 8 years ago. No FDIC consent obtained.}
\trap{Conviction over 5 years old; no consent required.}
\correct{Section 19 applies regardless of conviction age. FDIC consent required unless de minimis exception applies.}
\end{frame}

\begin{frame}{Mandatory Vacation - How It Detects Fraud}
\begin{block}{Employees Committing Fraud Avoid Vacation}
Mandatory vacation is a PROVEN fraud detection control.
\end{block}
\textbf{Why It Works:}
\begin{itemize}
\item Employee must take consecutive days off (5-10 business days)
\item Substitute performs duties during absence
\item Substitute may discover irregularities
\item Employee cannot conceal ongoing fraud while absent
\item Common for wire processors, account openers, override approvers
\end{itemize}
\textbf{Red Flags Related to Vacation:}
\begin{itemize}
\item Employee never takes vacation
\item Employee cancels vacation last minute repeatedly
\item Employee works during vacation
\item Employee refuses to provide passwords
\item Employee becomes anxious when vacation discussed
\end{itemize}
\redflag{Employees who never take vacation should be subject to enhanced scrutiny.}
\end{frame}

% ============================================================
% SECTION 12: QA AND ALERT ADJUDICATION (EXPANDED)
% ============================================================

\begin{frame}{QA Sample Size - Statistical Validity}
\begin{block}{QA Sample Must Be Representative}
Sample size must be sufficient to detect issues.
\end{block}
\textbf{Sampling Requirements:}
\begin{itemize}
\item Minimum 5-10 percent of total alerts (higher for high-risk areas)
\item Random sampling (not cherry-picked)
\item Targeted sampling (high-risk scenarios, specific analysts)
\item Sufficient volume to detect 5\% error rate with confidence
\item Document sampling methodology
\item Adjust sample size based on findings
\end{itemize}
\scenario{Bank reviews 2\% of 50,000 alerts annually. No issues found. QA program deemed effective.}
\trap{No issues found proves QA is effective.}
\correct{2\% sample may be insufficient to detect issues. Sample size must be statistically valid. Regulator may require larger sample.}
\end{frame}

\begin{frame}{Analyst Variation - Acceptable vs. Unacceptable Range}
\begin{block}{Some Variation Expected; Extreme Variation is Problem}
Controls needed when variation exceeds acceptable range.
\end{block}
\textbf{Acceptable Variation:}
\begin{itemize}
\item 5-10\% difference between analysts on similar caseloads
\item Based on legitimate judgment differences
\item Documented in case files
\end{itemize}
\textbf{Unacceptable Variation (Requires Action):}
\begin{itemize}
\item 30\%+ difference between analysts
\item Same scenario, different conclusions
\item One analyst closes 95\%, another escalates 70\%
\item Variation correlated with analyst tenure or workload
\end{itemize}
\redflag{Track analyst-level metrics. Wide variation indicates need for calibration.}
\end{frame}

% ============================================================
% SECTION 13: DATA QUALITY (EXPANDED)
% ============================================================

\begin{frame}{Data Governance - Who Owns Data Quality}
\begin{block}{Data Quality is Shared Responsibility}
Cannot be delegated solely to IT.
\end{block}
\textbf{Role Responsibilities:}
\begin{itemize}
\item Compliance: Define data requirements for monitoring
\item IT: Implement controls to meet requirements
\item Business: Ensure source system data quality
\item Data Governance: Establish enterprise standards
\item Internal Audit: Test data quality controls
\end{itemize}
\scenario{Data quality issues identified. IT says compliance didn't specify requirements. Compliance says IT should ensure data quality.}
\trap{IT is solely responsible for data quality.}
\correct{Compliance must define requirements. IT must implement controls. Shared responsibility.}
\end{frame}

\begin{frame}{Data Reconciliation - Daily Requirements}
\begin{block}{Must Reconcile Daily or Intra-Day}
Longer reconciliation cycles leave gaps undetected.
\end{block}
\textbf{Reconciliation Controls:}
\begin{itemize}
\item Source system transaction count vs. monitoring system count (daily)
\item Source system dollar volume vs. monitoring system volume (daily)
\item Hash total comparisons for data integrity
\item Exception report for any discrepancies
\item Resolution within 24-48 hours
\item Escalation for unresolved exceptions
\item Board reporting for material gaps
\end{itemize}
\redflag{Daily reconciliation is minimum expectation for most institutions.}
\end{frame}

% ============================================================
% SECTION 14: NEW PRODUCT RISK ASSESSMENT (EXPANDED)
% ============================================================

\begin{frame}{New Product Risk Assessment - Pre-Launch Requirements}
\begin{block}{FATF Recommendation 15 - New Technologies}
Risk assessment must be completed BEFORE product launch.
\end{block}
\textbf{Required Assessment Components:}
\begin{itemize}
\item ML/TF vulnerability assessment
\item Anonymity features analysis
\item Cross-border implications
\item Third-party payment capabilities
\item Transaction speed and volume potential
\item Regulatory compliance review
\item Control gap identification
\item Remediation plan for identified gaps
\end{itemize}
\scenario{Bank launches new mobile payment product. AML risk assessment completed 3 months AFTER launch. No issues identified.}
\trap{Risk assessment after launch is acceptable if no issues found.}
\correct{FATF Rec. 15 requires risk assessment BEFORE launch. Post-launch assessment is violation regardless of findings.}
\end{frame}

\begin{frame}{New Product - Compliance Officer Involvement}
\begin{block}{CCO Must Be Active Participant}
Compliance cannot be consulted after product design is complete.
\end{block}
\textbf{Required CCO Involvement:}
\begin{itemize}
\item Initial product concept review
\item Design phase participation
\item Control requirement definition
\item Risk assessment approval
\item Sign-off before launch
\item Post-launch monitoring review
\end{itemize}
\scenario{Product team designs new payment product. Presents completed product to CCO for sign-off. CCO identifies AML control gaps requiring redesign.}
\trap{CCO review at completion is sufficient.}
\correct{CCO must be involved FROM THE START. Post-design review is insufficient and causes delays.}
\end{frame}

% ============================================================
% SECTION 15: CONCLUSION - ALL MISSING CONCEPTS SUMMARY
% ============================================================

\begin{frame}{Domain C Missing Content - Complete Summary}
\begin{multicols}{2}
\begin{enumerate}
\item \textbf{Risk Assessment:} Inherent vs. residual calculation
\item \textbf{Risk Appetite:} Board-approved thresholds
\item \textbf{Dynamic Reassessment:} Trigger events
\item \textbf{Board Minutes:} Required discussion content
\item \textbf{CCO Independence:} Reporting line requirements
\item \textbf{CCO Resources:} Adequate staffing mandates
\item \textbf{SAR Deadlines:} 30/60 day rule
\item \textbf{CTR Deadline:} 15 day rule
\item \textbf{314(a) Deadline:} 14 business days
\item \textbf{FBAR Deadline:} April 15 / October 15
\item \textbf{Section 311:} Three levels of measures
\item \textbf{Section 313:} Shell bank prohibition
\item \textbf{Section 319(b):} Record production
\item \textbf{314(a) vs. 314(b):} Key differences
\item \textbf{314(b) Opt-In:} Required before sharing
\item \textbf{Vendor Due Diligence:} Pre-engagement requirements
\item \textbf{Right-to-Audit:} Contract requirement
\item \textbf{Independent Testing:} Operational vs. design
\item \textbf{Legal Hold:} Subpoena triggers preservation
\item \textbf{Recordkeeping:} 5 years from specific dates
\item \textbf{Escalation Matrix:} Who approves what
\item \textbf{Keep Open Requests:} Written documentation required
\item \textbf{Section 19:} Disqualified persons
\item \textbf{Mandatory Vacation:} Fraud detection control
\item \textbf{QA Sample Size:} Statistical validity
\item \textbf{Data Governance:} Shared responsibility
\item \textbf{New Product Assessment:} Pre-launch requirement
\item \textbf{CCO Product Involvement:} From start, not end
\end{enumerate}
\end{multicols}
\greennote{These 28 concepts are consistently reported as surprise questions on recent CAMS exams.}
\end{frame}

\begin{frame}{Critical Red Flags - Domain C Missing Content}
\begin{itemize}
\item Board minutes with one-line approval only → Insufficient oversight
\item CCO reports to Head of Sales → Independence violation
\item \$100M client delays SAR → Commercial pressure violation
\item 40,000 alert backlog → Inadequate resources
\item SAR filed on Day 75 with suspect known at Day 10 → Missed deadline
\item 314(b) sharing without opt-in → Tipping-off violation
\item Vendor due diligence AFTER contract → Pre-engagement failure
\item Audit tests policies only, not transactions → Design-only testing
\item Records deleted 5 years from account opening → Premature deletion
\item Subpoena received, records deleted next day → Spoliation
\item FBI oral keep open request, no documentation → Insufficient
\item Funds withheld after account closure → Conversion
\item Fraud convict hired without FDIC consent → Section 19 violation
\item Employee never takes vacation → Insider risk red flag
\item 2\% QA sample, no issues found → Statistically invalid
\item New product launched, risk assessment 3 months later → Rec. 15 violation
\item CCO consulted at product completion → Too late for controls
\end{itemize}
\end{frame}

\begin{frame}{Exam Day Reminders - Domain C Missing Questions}
\begin{itemize}
\item \textbf{Inherent vs. Residual:} Residual = Inherent ÷ Control effectiveness factor
\item \textbf{Risk Appetite:} Board must define acceptable residual risk levels
\item \textbf{Reassessment Triggers:} SAR filing, negative media, ownership change
\item \textbf{Board Minutes:} Must show discussion, not just approval
\item \textbf{CCO Independence:} Cannot report to revenue function
\item \textbf{SAR Deadlines:} 30 days (suspect known) / 60 days (suspect unknown)
\item \textbf{CTR Deadline:} 15 days from transaction
\item \textbf{314(a):} MANDATORY, 14 BUSINESS days
\item \textbf{314(b):} Voluntary, requires OPT-IN before sharing
\item \textbf{Section 311:} Applies to FOREIGN institutions
\item \textbf{Section 313:} Shell banks ABSOLUTELY PROHIBITED
\item \textbf{Vendor Due Diligence:} BEFORE contract, not after
\item \textbf{Right-to-Audit:} Must be in contract, enforceable
\item \textbf{Independent Testing:} Must test OPERATIONAL effectiveness
\item \textbf{Legal Hold:} Automatic upon subpoena receipt
\item \textbf{Record Retention:} 5 years from specific dates (see table)
\item \textbf{Keep Open Requests:} Must be IN WRITING
\item \textbf{Section 19:} FDIC consent required for dishonesty convictions
\item \textbf{Mandatory Vacation:} Detect fraud, identify red flags
\item \textbf{New Products:} Risk assessment BEFORE launch
\item \textbf{CCO Involvement:} From product conception, not after design
\end{itemize}
\vspace{0.3cm}
\centering
\greennote{Good luck on your CAMS exam - you've covered the hard Domain C questions now!}
\end{frame}

\end{document}